\documentclass[11pt]{article}
% BEGIN SHARED COURSE STYLE
% Shared Math 615 style. Embedded verbatim in standalone published sources.
\RequirePackage[letterpaper,margin=1in]{geometry}
\RequirePackage{amsmath,amsthm,mathtools}
\RequirePackage{fontspec,unicode-math}
\setmainfont{texgyretermes}[Extension=.otf,UprightFont=*-regular,BoldFont=*-bold,ItalicFont=*-italic,BoldItalicFont=*-bolditalic]
\setmathfont{texgyretermes-math.otf}
\RequirePackage{microtype,tikz-cd,enumitem,needspace,xcolor,booktabs,titlesec,etoolbox}
\definecolor{teal}{HTML}{176568}
\definecolor{burgundy}{HTML}{782F40}
\definecolor{inklink}{HTML}{176568}
\definecolor{accent}{HTML}{176568}
\titleformat{\section}{\large\bfseries\color{teal}}{\thesection}{.65em}{}
\titleformat{\subsection}{\normalsize\bfseries\color{burgundy}}{\thesubsection}{.65em}{}
\titleformat{\paragraph}[runin]{\normalsize\bfseries\color{burgundy}}{}{0pt}{}
\titlespacing*{\section}{0pt}{2.2ex plus .5ex minus .2ex}{1ex plus .2ex}
\titlespacing*{\subsection}{0pt}{1.6ex plus .4ex minus .2ex}{.7ex plus .2ex}
\RequirePackage{hyperref,bookmark}
\hypersetup{colorlinks=true,linkcolor=teal,urlcolor=teal,citecolor=teal,pdfstartview={XYZ null null 1.5},bookmarksnumbered=true}
\urlstyle{same}
\setcounter{tocdepth}{2}
\setlist[enumerate]{label=\arabic*.,ref=\arabic*,leftmargin=*,itemsep=4pt,topsep=4pt}
\setlength{\parindent}{1em}
\setlength{\parskip}{3pt}
\setlength{\emergencystretch}{2em}
\raggedbottom
\AddToHook{cmd/section/before}{\Needspace{8\baselineskip}}
\AddToHook{cmd/subsection/before}{\Needspace{6\baselineskip}}
\makeatletter
\newcommand{\afterblock}{\@afterindentfalse\@afterheading}
\makeatother
\newcounter{result}[section]
\renewcommand{\theresult}{\thesection.\arabic{result}}
\newcommand{\resulttitle}[3]{#1 #2\ifstrempty{#3}{}{ (#3)}.}
\newenvironment{result}[3]{\par\addvspace{7pt}\Needspace{4\baselineskip}\refstepcounter{result}\hypertarget{#2}{}\label{#2}\noindent{\color{burgundy}\hypersetup{linkcolor=burgundy}\hyperlink{proof:#2}{\textbf{\resulttitle{#1}{\theresult}{#3}}}}\quad\ignorespaces}{\par\addvspace{4pt}\afterblock}
\newenvironment{entry}[3]{\par\addvspace{7pt}\Needspace{3\baselineskip}\refstepcounter{result}\hypertarget{#2}{}\label{#2}\noindent{\color{burgundy}\textbf{\resulttitle{#1}{\theresult}{#3}}}\quad\ignorespaces}{\par\addvspace{4pt}\afterblock}
\newenvironment{demonstration}[1]{\par\noindent\hypertarget{proof:#1}{}{\color{burgundy}\hypersetup{linkcolor=burgundy}\hyperlink{#1}{\textbf{Proof.}}}\quad\ignorespaces}{\unskip\nobreak\hfill\qedsymbol\par\addvspace{5pt}\afterblock}
\newenvironment{citedresult}[4]{\par\addvspace{7pt}\Needspace{4\baselineskip}\refstepcounter{result}\hypertarget{#2}{}\label{#2}\noindent{\color{burgundy}\hypersetup{urlcolor=burgundy}\href{#4}{\textbf{\resulttitle{#1}{\theresult}{#3}}}}\quad\ignorespaces}{\par\addvspace{4pt}\afterblock}
\newcommand{\usedby}[1]{\par\noindent{\small\textbf{Used in:} #1.}\par\afterblock}
\newcommand{\coursebold}[1]{{\color{burgundy}\textbf{#1}}}
\newcommand{\coursetitle}[3]{\begin{center}{\large\bfseries\color{teal}#1}\\[4pt]{\LARGE\bfseries\color{teal}#2}\\[6pt]Math 615: Algebraic Topology I\\Fall 2026\quad\textbar\quad #3\end{center}}
\newcommand{\coursecontents}{{\small\setlength{\parskip}{0pt}\tableofcontents}}
\newcommand{\homeworkstyle}{\titleformat{\section}{\large\bfseries\color{teal}}{Exercise \thesection.}{.65em}{}}
% END SHARED COURSE STYLE
\newcommand{\Mor}{\operatorname{Mor}}
\newcommand{\Tor}[3]{\hyperlink{def:tor}{\operatorname{Tor}^{R}_{#1}(#2,#3)}}
\newcommand{\Extn}[3]{\hyperlink{def:ext}{\operatorname{Ext}^{#1}_{R}(#2,#3)}}
\newcommand{\Hc}[2]{\hyperlink{def:cochain}{H^{#1}(#2)}}
\newcommand{\im}{\operatorname{im}}
\newcommand{\coker}{\operatorname{coker}}
\newcommand{\Ext}{\operatorname{Ext}}
\newcommand{\id}{\mathrm{id}}
\newcommand{\A}{\mathcal A}
\newcommand{\Z}{\mathbb Z}
\newcommand{\Ch}[1]{\hyperlink{def:complex}{\operatorname{Ch}(#1)}}
\newcommand{\K}[1]{\hyperlink{def:K}{K(#1)}}
\newcommand{\D}[1]{\hyperlink{def:D}{D(#1)}}
\newcommand{\Hn}[2]{\hyperlink{def:homology}{H_{#1}(#2)}}
\newcommand{\Zn}[2]{\hyperlink{def:homology}{Z_{#1}(#2)}}
\newcommand{\Bn}[2]{\hyperlink{def:homology}{B_{#1}(#2)}}
\newcommand{\shift}[2]{\hyperlink{def:shift}{#1[#2]}}
\newcommand{\cone}[1]{\hyperlink{def:cone}{\operatorname{Cone}(#1)}}
\newcommand{\cyl}[1]{\hyperlink{def:cylinder}{\operatorname{Cyl}(#1)}}

\DeclareMathOperator{\Cone}{Cone}
\DeclareMathOperator{\Cyl}{Cyl}



\hypersetup{pdftitle={Lecture 2: Exact Sequences and Derived Categories},pdfauthor={Math 615 Fall 2026},pdfsubject={Revision 6}}
\begin{document}
\coursetitle{Lecture 2}{Exact Sequences and Derived Categories}{September 9}
\noindent
The long exact sequence, its mapping-cone interpretation, and the passage to $K(\A)$ and $D(\A)$. Section~\ref{sec:size} is supplementary reading.

\coursecontents

\section{From a Short Exact Sequence to Homology}\label{sec:les}
\subsection{Conventions and the Connecting Morphism}
Let $\A$ be a locally small abelian category. We use homological grading. Recall:
\hypertarget{def:complex}{}a complex has differentials $d_n:C_n\to C_{n-1}$ with $d^2=0$; chain maps commute with differentials. Write $\Ch{\A}$ for their category.
\hypertarget{def:homology}{}Set $Z_n(C)=\ker d_n$, $B_n(C)=\im d_{n+1}$, and
\[
H_n(C)=\coker\bigl(B_n(C)\longrightarrow Z_n(C)\bigr).
\]
A complex is \emph{acyclic} if its homology is zero. A chain map is a \emph{quasi-isomorphism} if every induced homology map is an isomorphism.

\subsection{The Snake Lemma}
\begin{result}{Lemma}{lem:snake}{Kernels, Cokernels, and the Connecting Map}
Consider a commutative diagram in an abelian category with exact rows:
\[
\begin{tikzcd}[column sep=large,row sep=large]
&A\arrow[r,"u"]\arrow[d,"\alpha"']&B\arrow[r,"v"]\arrow[d,"\beta"]&C\arrow[r]\arrow[d,"\gamma"]&0\\
0\arrow[r]&A'\arrow[r,"u'"']&B'\arrow[r,"v'"']&C'.&
\end{tikzcd}
\]
There is a natural exact sequence
\[
\ker\alpha\longrightarrow\ker\beta\longrightarrow\ker\gamma
\xrightarrow{\delta}\coker\alpha\longrightarrow\coker\beta
\longrightarrow\coker\gamma.
\]
If $u$ is monic, the sequence starts with $0$; if $v'$ is epic, it ends with $0$.
\end{result}
\begin{demonstration}{lem:snake}
Let $K=\ker\gamma$ and form $T=B\times_C K$. Its projection $p:T\to K$
is epic because it is a pullback of the epimorphism $v$. The composite
$T\to B\xrightarrow{\beta}B'$ is killed by $v'$, so it factors uniquely
as $u't$ for $t:T\to A'$. The induced map $A\to T$ has image $\ker p$,
and $t$ composed with this map is $\alpha$. Thus $q_\alpha t$ vanishes
on $\ker p$, where $q_\alpha:A'\to\coker\alpha$ is the quotient.
There is a unique $\delta:K\to\coker\alpha$ satisfying
$\delta p=q_\alpha t$.

For clarity, the following lift calculations may be performed after
pullback along epimorphisms. A lift of a map through an epimorphism exists
on its pullback; such local lifts suffice to compare images and kernels,
because epimorphisms detect equality of subobjects. Thus the calculations
are valid in an abelian category and require no elements or chosen sections.
At $\ker\beta$, a map whose composite with $v$ vanishes locally has the
form $ua$; the equation $u'\alpha a=\beta ua=0$ gives $\alpha a=0$.
At $\ker\gamma$, lift $c$ to $b$ and write $\beta b=u'a'$.
If $\delta c=0$, lift $a'$ through the image of $\alpha$, writing
$a'=\alpha a$ locally. Then $b-ua$ is killed by $\beta$ and still maps to $c$.
Conversely a lift killed by $\beta$ gives zero connecting class.

At $\coker\alpha$, a representative $a'$ maps to zero in $\coker\beta$
exactly when $u'a'=\beta b$ locally. Then $vb$ is killed by $\gamma$
and its connecting class is $[a']$. At $\coker\beta$, suppose a
representative $b'$ maps to zero in $\coker\gamma$. Locally write
$v'b'=\gamma c$ and lift $c=vb$. Then $b'-\beta b$ is killed by $v'$
and hence lies in the image of $u'$, as required. The reverse inclusions
follow from commutativity. These arguments prove exactness at all four
interior terms. The asserted endpoint injectivity and surjectivity follow
from those of $u$ and $v'$. The pullback construction commutes with maps
of the original diagrams, proving naturality.
\end{demonstration}
\usedby{Theorem~\ref{thm:les}}

For modules the construction is the familiar rule: lift $c\in\ker\gamma$
to $b\in B$, solve $u'a'=\beta b$, and put $\delta(c)=[a']$.
Changing $b$ to $b+ua$ changes $a'$ to $a'+\alpha a$, so the cokernel
class does not change.

\begin{result}{Theorem}{thm:les}{The Long Exact Sequence}
A degreewise short exact sequence of complexes
\begin{equation}\label{eq:ses}
0\longrightarrow A\xrightarrow{i}B\xrightarrow{\rho}C\longrightarrow0
\end{equation}
induces natural connecting morphisms $\delta_n:\Hn{n}{C}\to\Hn{n-1}{A}$ and an exact sequence
\begin{equation}\label{eq:les}
\cdots\longrightarrow\Hn{n}{A}\xrightarrow{\Hn{n}{i}}\Hn{n}{B}
\xrightarrow{\Hn{n}{\rho}}\Hn{n}{C}\xrightarrow{\delta_n}\Hn{n-1}{A}
\longrightarrow\cdots.
\end{equation}
For modules, if $c\in C_n$ is a cycle, choose $b\in B_n$ with $\rho b=c$ and write $d^B b=ia$. Then $a$ is a cycle and $\delta_n$ sends the class of $c$ to the class of $a$.
\end{result}

\begin{demonstration}{thm:les}
Put $Q_n(X)=\coker(d^X_{n+1})=X_n/\Bn{n}{X}$. The differential induces $\partial_X:Q_n(X)\to\Zn{n-1}{X}$. Its kernel is $\Hn{n}{X}$ and its cokernel is $\Hn{n-1}{X}$: its image is precisely $\Bn{n-1}{X}$.

The following commutative diagram has exact rows:
\begin{equation}\label{eq:snake}
\begin{tikzcd}[column sep=large,row sep=large]
&Q_n(A)\arrow[r]\arrow[d,"\partial_A"']&Q_n(B)\arrow[r]\arrow[d,"\partial_B"]&Q_n(C)\arrow[r]\arrow[d,"\partial_C"]&0\\
0\arrow[r]&\Zn{n-1}{A}\arrow[r]&\Zn{n-1}{B}\arrow[r]&\Zn{n-1}{C}.&
\end{tikzcd}
\end{equation}
The upper row is exact because cokernels preserve right exact sequences of arrows, applied to the differentials in degrees $n+1,n$. The lower row is exact because kernels preserve left exact sequences of arrows, applied in degrees $n-1,n-2$. 

Lemma~\ref{lem:snake}, in its form with a right exact upper row and a left exact lower row, gives
\[
\ker\partial_A\to\ker\partial_B\to\ker\partial_C
\xrightarrow{\delta_n}\coker\partial_A\to\coker\partial_B\to\coker\partial_C.
\]
Substituting the identifications above proves exactness at all three types of homology terms. The six-term sequences for consecutive $n$ join because their other maps are the maps induced by $i,\rho$. Naturality follows from the snake lemma; see \cite[Tag 0117]{Sta26}.

For modules, $\rho d^B b=d^C c=0$, so exactness gives a unique $a$ with $ia=d^B b$. Since $i(d^A a)=d^B d^B b=0$ and $i$ is mono, $d^A a=0$. Changing $b$ to $b+ix$ changes $a$ to $a+d^A x$. Changing $c$ to $c+d^C c'$ and choosing a lift $b'$ of $c'$ changes a chosen lift to $b+d^B b'$ and leaves its differential unchanged. Thus the formula is well defined and agrees with the snake map. 
\end{demonstration}
\usedby{Sections~\ref{sec:cone} and~\ref{sec:derived}}

\begin{entry}{Example}{ex:relative}{Relative Homology}
If $A\subseteq B$ is a subcomplex, define the quotient complex $B/A$ degreewise. The theorem gives
\[
\cdots\to\Hn{n}{A}\to\Hn{n}{B}\to\Hn{n}{B/A}\xrightarrow{\delta_n}\Hn{n-1}{A}\to\cdots.
\]
For singular chains of a pair of spaces $(X,Y)$, take $B=C_*(X;R)$ and $A=C_*(Y;R)$. A relative cycle is represented by a chain in $X$ whose boundary lies in $Y$; $\delta$ takes the homology class of that boundary in $Y$.
\end{entry}

\section{Mapping Cones and Applications}\label{sec:cone}
\subsection{The Cone Sequence}
\begin{entry}{Definition}{def:shift}{Shifts}
For $r\in\Z$, set $C[r]_n=C_{n-r}$ and $d^{C[r]}_n=(-1)^r d^C_{n-r}$. A chain map shifts by $f[r]_n=f_{n-r}$. Thus $\Hn{n}{\shift{C}{r}}=\Hn{n-r}{C}$ canonically. For any biproduct, $\iota_j$ and $\rho_j$ denote the canonical inclusions and projections.
\end{entry}

\begin{entry}{Definition}{def:cone}{The Mapping Cone}
For a chain map $f:C\to D$, define
\[
\cone{f}_n=D_n\oplus C_{n-1},\qquad
 d^{\cone{f}}_n=\begin{pmatrix}d^D_n&f_{n-1}\\0&-d^C_{n-1}\end{pmatrix}.
\]
The upper-right entry of its squared differential is $d^D_{n-1}f_{n-1}-f_{n-2}d^C_{n-1}=0$; all other entries vanish. There is a degreewise split short exact sequence
\begin{equation}\label{eq:cone-ses}
0\to D\xrightarrow{j=\iota_1}\cone{f}\xrightarrow{q=\rho_2}\shift{C}{1}\to0.
\end{equation}
\end{entry}

\begin{result}{Proposition}{prop:cone-les}{Homology of a Cone}
The connecting map of \eqref{eq:cone-ses} is $\Hn{n-1}{f}$. Hence
\begin{equation}\label{eq:cone-les}
\cdots\to\Hn{n}{C}\xrightarrow{\Hn{n}{f}}\Hn{n}{D}\xrightarrow{\Hn{n}{j}}\Hn{n}{\cone{f}}
\xrightarrow{\Hn{n}{q}}\Hn{n-1}{C}\xrightarrow{\Hn{n-1}{f}}\Hn{n-1}{D}\to\cdots
\end{equation}
is exact. In particular, $f$ is a quasi-isomorphism if and only if $\cone{f}$ is acyclic.
\end{result}
\begin{demonstration}{prop:cone-les}
A cycle $c\in C_{n-1}$ lifts to $(0,c)$, whose cone differential is $(fc,0)$. Thus the connecting map has the asserted positive sign. Categorically, the second-summand section of \eqref{eq:cone-ses} has the same differential defect $f$, so this computes the snake map in $\A$. Apply Theorem~\ref{thm:les}. If every $\Hn{n}{f}$ is invertible, exactness forces all cone homology to vanish. Conversely, its vanishing makes $\Hn{n}{f}$ both mono and epi, hence an isomorphism in $\A$.
\end{demonstration}
\usedby{Proposition~\ref{prop:ses-cone}}

\begin{entry}{Example}{ex:multiplication}{Multiplication by an Integer}
For $C=D=\Z$ concentrated in degree zero and $f=m$, the cone is $\Z\xrightarrow{m}\Z$ in degrees $1,0$. Its homology is $\ker(m)$ in degree $1$ and $\Z/m\Z$ in degree $0$.

\end{entry}

\subsection{Recovering a Short Exact Sequence from a Cone}
\begin{result}{Proposition}{prop:ses-cone}{A Cone Models the Quotient}
For \eqref{eq:ses}, let $\pi:\cone{i}\to C$ have restrictions $\rho:B\to C$ and $0:\shift{A}{1}\to C$. Then $\pi$ is a quasi-isomorphism and the diagram
\[
\begin{tikzcd}[column sep=large,row sep=small]
A\arrow[r,"i"]\arrow[d,equal]&B\arrow[r,"j"]\arrow[d,equal]&\cone{i}\arrow[d,"\pi","\sim"']\\
A\arrow[r,"i"']&B\arrow[r,"\rho"']&C
\end{tikzcd}
\]
commutes. With $q:\cone{i}\to\shift{A}{1}$ the second projection, the connecting morphism in Theorem~\ref{thm:les} is
\begin{equation}\label{eq:delta-cone}
\delta_n=-\Hn{n}{q}\,\Hn{n}{\pi}^{-1}.
\end{equation}
\end{result}
\begin{demonstration}{prop:ses-cone}
The map $\pi$ is a degreewise epimorphism of complexes. Its kernel identifies with $\cone{\id_A}$: the inclusion sends $(x,a)$ to $(ix,a)$. On this kernel, $s_n(x,a)=(0,x)$ is a contracting homotopy, since $ds+sd=\id$. Theorem~\ref{thm:les} therefore implies that $\pi$ is a quasi-isomorphism.

For the sign, lift a cycle $c$ to $b$ and write $db=ia$. The corresponding cone cycle is $(b,-a)$, so $q(b,-a)=-a$, while $\delta_n$ sends the class of $c$ to the class of $a$. The same computation with the biproduct maps gives the identity in $\A$.
\end{demonstration}
\usedby{Definition~\ref{def:triangles}}
Thus the cone sequence recovers the long exact sequence of \eqref{eq:ses}; formula~\eqref{eq:delta-cone} identifies the chosen connecting map, including its sign.

\section{The Homotopy and Derived Categories}\label{sec:derived}
\subsection{Passing to Chain Homotopy Classes}
\hypertarget{def:homotopy}{}A chain homotopy from $g$ to $f$ is a family $h_n:C_n\to D_{n+1}$ with $f-g=d^Dh+hd^C$. A complex is \emph{contractible} if its identity is nullhomotopic.

\begin{entry}{Definition}{def:K}{The Homotopy Category}
Let $N(C,D)$ be the subgroup of nullhomotopic chain maps. The category $K(\A)$ has the same objects as $\Ch{\A}$ and morphism groups
\[
\Mor_{K(\A)}(C,D)=\Mor_{\Ch{\A}}(C,D)/N(C,D).
\]
Composition is induced by composition of chain maps; homotopy is as above.
\end{entry}
\begin{result}{Proposition}{prop:homotopy}{Well-Defined Composition and Homology}
Nullhomotopic maps form an additive ideal. Hence $K(\A)$ is additive and every homology functor factors through it. Its isomorphisms are precisely chain homotopy equivalences.
\end{result}
\begin{demonstration}{prop:homotopy}
The sum of homotopies is a homotopy for the sum. If $f=dh+hd$ and $a,b$ are composable chain maps, then $bfa=d(bha)+(bha)d$. Thus composition descends to the quotient. On cycles, $f-g=dh$, whose image is a boundary; therefore $f,g$ induce the same homology map. An inverse in the quotient means exactly inverse chain maps up to homotopy.
\end{demonstration}
\usedby{Proposition~\ref{prop:small-K}}

\begin{entry}{Example}{ex:acyclic}{Acyclic Is Weaker Than Contractible}
The exact complex $0\to\Z\xrightarrow{2}\Z\to\Z/2\to0$ in degrees $2,1,0$ is acyclic. A contracting homotopy would split the quotient $\Z\to\Z/2$ in degree zero, which is impossible. It is therefore nonzero in $K(\mathbf{Ab})$.
\end{entry}

\subsection{Inverting Quasi-Isomorphisms}
Assume here that $\A$ is essentially small or Grothendieck (for example, a category of modules). The size issue is addressed in Section~\ref{sec:size}.

\begin{entry}{Definition}{def:D}{The Derived Category}
Let $S$ be the class of quasi-isomorphisms in $K(\A)$. Define
\[
D(\A)=K(\A)[S^{-1}].
\]
It has the same objects as $K(\A)$ and a functor $Q:K(\A)\to D(\A)$ sending each $s\in S$ to an isomorphism. Its universal property is that precomposition with $Q$ identifies functors out of $D(\A)$ with functors out of $K(\A)$ which invert $S$; natural transformations correspond as well.
\end{entry}

Morphisms are represented by roofs in $K(\A)$:
\[
\begin{tikzcd}[column sep=large,row sep=small]
&Z\arrow[dl,"s"',"\sim"]\arrow[dr,"u"]&\\
C&&D.
\end{tikzcd}
\qquad\text{The represented map is }Q(u)Q(s)^{-1}.
\]
Two roofs are equivalent if they admit a common refinement in $K(\A)$ on which both maps agree, with the common map to $C$ a quasi-isomorphism. This roof description is the calculus of fractions for quasi-isomorphisms; its construction is not needed here.

Homology descends to $D(\A)$ by the universal property. A complex becomes zero precisely when it is acyclic: use the quasi-isomorphism $0\to C$ in one direction, and homology in the other. Similarly, a chain map becomes invertible precisely when it is a quasi-isomorphism.

\begin{entry}{Definition}{def:triangles}{The Connecting Map in the Derived Category}
The quasi-isomorphism $\pi:\cone{i}\to C$ gives a morphism
\[
\partial=Q(-q)Q(\pi)^{-1}:C\longrightarrow\shift{A}{1}
\]
in $D(\A)$, with $\Hn{n}{\partial}=\delta_n$. Thus the short exact sequence determines
\[
A\xrightarrow{i}B\xrightarrow{\rho}C\xrightarrow{\partial}\shift{A}{1}.
\]
This is its distinguished triangle. With our signs, cone triangles have the form
$X\xrightarrow{f}Y\xrightarrow{j}\cone{f}\xrightarrow{-q}\shift{X}{1}$.
\end{entry}


\section{Local Smallness: A Final Supplement}\label{sec:size}
\begin{result}{Proposition}{prop:small-K}{Local Smallness of the Homotopy Category}
If $\A$ is locally small, then $\Ch{\A}$ and $K(\A)$ are locally small, even for unbounded complexes.
\end{result}
\begin{demonstration}{prop:small-K}
For fixed $C,D$, every chain map lies in the set
\[
\prod_{n\in\Z}\Mor_{\A}(C_n,D_n).
\]
The chain-map equations select a subset. The nullhomotopic maps form a subgroup of this set-sized abelian group; quotienting by it gives a set. The countable product here is a product of sets, not a product that must exist in $\A$.
\end{demonstration}

\begin{result}{Proposition}{prop:small-D}{Essentially Small Abelian Categories}
If $\A$ is essentially small, then $D(\A)$ is locally small.
\end{result}
\begin{demonstration}{prop:small-D}
Replace $\A$ by a small skeleton. Its complexes and chain maps form sets. Finite words in arrows and formal inverses also form a set; localization quotients this set by relations.
\end{demonstration}

\begin{result}{Proposition}{prop:small-mod}{A Direct Size Bound for Modules}
For a set-sized ring $R$, the unbounded category $D(R\text{-}\mathbf{Mod})$ is locally small.
\end{result}
\begin{demonstration}{prop:small-mod}
Fix complexes $C,D$, and an infinite cardinal $\kappa$ at least $|R|$, $\aleph_0$, and the cardinality of the disjoint union of the underlying sets $C_n$. Consider any roof $C\xleftarrow{s}E\xrightarrow{u}D$ with $s$ a quasi-isomorphism.

Choose cycles in $E$ representing every homology class of $C$. They generate a subcomplex $E^{(0)}$ of total cardinality at most $\kappa$. From $E^{(m)}$, form $E^{(m+1)}$ by adjoining one bounding element in $E$ for each cycle whose image bounds in $C$, then closing under the differential and module operations. These bounds exist since $H_*(s)$ is injective; the cardinality remains at most $\kappa$.

The union $E'=\bigcup_{m\ge0}E^{(m)}$ still has size at most $\kappa$, and $s|_{E'}$ is a quasi-isomorphism: chosen cycles give surjectivity, and the successive bounds give injectivity. Thus every roof refines to one whose middle complex has size at most $\kappa$. There is a set of such complexes up to isomorphism and a set of maps from each to $C,D$. Quotienting this set of roofs gives the required morphism set.
\end{demonstration}

\begin{entry}{Remark}{rem:size}{The General Case}
A Grothendieck abelian category has small coproducts, exact filtered colimits, and a generator. Its derived category is locally small by the existence of K-injective resolutions \cite[Tags 070G, 079P]{Sta26}; that theorem is beyond this lecture. For an arbitrary locally small abelian category the conclusion can fail, even for the bounded derived category \cite{CN08}.
\end{entry}

\begin{thebibliography}{Sta26}
\raggedright
\bibitem[Sta26]{Sta26} The Stacks Project Authors, \emph{The Stacks Project}, 2026. \href{https://stacks.math.columbia.edu/tag/0117}{Tag 0117: Long Exact Cohomology Sequence}; \href{https://stacks.math.columbia.edu/tag/014D}{Tag 014D: Cones and Termwise Split Sequences}; \href{https://stacks.math.columbia.edu/tag/070G}{Tag 070G: K-Injective Complexes}; \href{https://stacks.math.columbia.edu/tag/079P}{Tag 079P: Existence of K-Injective Resolutions}. The source uses cohomological grading; these notes use homological grading.
\bibitem[CN08]{CN08} Carles Casacuberta and Amnon Neeman, \emph{Brown Representability Does Not Come for Free}, arXiv:0807.1872v1, 2008. \href{https://arxiv.org/abs/0807.1872v1}{arXiv}. See Section 1 for the failure of small morphism sets.
\end{thebibliography}
\end{document}
